\SourceLeaf{250}%
ficiei ductum, atque \(q\) angulum, quem primum elementum huius lineae efficit cum elemento primo alicuius lineae brevissimae ex \(A\) proficiscentis datae. Sit \(B\) punctum determinatum in hac linea pro qua \(q=0\), atque \(C\) aliud punctum determinatum superficiei, pro quo valorem ipsius \(q\) simpliciter per \(A\) designabimus. Supponamus, puncta \(B,C\) per lineam brevissimam iuncta, cuius partes, inde a puncto \(B\) numeratas, indefinite ut in art. 18 per \(s\) denotabimus, nec non perinde ut illic, per \(\theta\) angulum, quem quodvis elementum \(\mathrm{d}s\) facit cum elemento \(\mathrm{d}p\): denique sint \(\theta^0,\theta'\) valores anguli \(\theta\) in punctis \(B,C\). Habemus itaque in superficie curva triangulum lineis brevissimis inclusum, eiusque anguli ad \(B\) et \(C\), per has ipsas literas simpliciter designandi aequales erunt ille complemento anguli \(\theta^0\) ad \(180^0\), hic ipsi angulo \(\theta'\). Sed quum analysin nostram inspicienti facile pateat, omnes angulos non per gradus sed per numeros expressos concipi, ita ut angulus \(57^0\,17'\,45''\), cui respondet arcus radio aequalis, pro unitate habeatur, statuere oportet, denotando per \(2\pi\) peripheriam circuli
\[
\theta^0=\pi-B,\qquad\theta'=C
\]
Inquiramus nunc in curvaturam integram huius trianguli, quae fit \(=\int k\mathrm{d}\sigma\), denotante \(\mathrm{d}\sigma\) elementum superficiale trianguli; quare quum hoc elementum exprimatur per \(m\mathrm{d}p\mathord{.}\mathrm{d}q\), eruere oportet integrale \(\int\!\int km\mathrm{d}p\mathord{.}\mathrm{d}q\) supra totam trianguli superficiem. Incipiamus ab integratione secundum \(p\), quae propter \(k=-\frac{1}{m}\mathord{.}\frac{\mathrm{d}\mathrm{d}m}{\mathrm{d}p^2}\), suppeditat \(\mathrm{d}q\mathord{.}(\mathrm{Const.}-\frac{\mathrm{d}m}{\mathrm{d}p})\) pro curvatura integra areae iacentis inter lineas primi systematis, quibus respondent valores indeterminatae secundae \(q,q+\mathrm{d}q\): quum haec curvatura pro \(p=0\) evanescere debeat, quantitas constans per integrationem introducta aequalis esse debet valori ipsius \(\frac{\mathrm{d}m}{\mathrm{d}p}\) pro \(p=0\), i. e. unitati. Habemus itaque \(\mathrm{d}q(1-\frac{\mathrm{d}m}{\mathrm{d}p})\), ubi pro \(\frac{\mathrm{d}m}{\mathrm{d}p}\) accipere oportet valorem respondentem fini illius areae in linea \(CB\). In hac linea vero fit per art. praec. \(\frac{\mathrm{d}m}{\mathrm{d}p}\mathord{.}\mathrm{d}q=-\mathrm{d}\theta\), unde expressio nostra mutatur in \(\mathrm{d}q+\mathrm{d}\theta\). Accedente iam integratione altera a \(q=0\) usque ad \(q=A\) extendenda, obtinemus curvaturam integram trianguli \(=A+\theta'-\theta^0=A+B+C-\pi\).

Curvatura integra aequalis est areae eius partis superficiei sphaericae, quae respondet triangulo, signo positivo vel negativo affectae, prout superficies curva, in qua triangulum iacet, est concavo-concava vel concavo-convexa: pro unitate areae accipiendum est quadratum, cuius latus est unitas (radius sphaerae), quo pacto superficies tota sphaerae fit \(=4\pi\). Est itaque pars superficiei sphaericae
